\textbf{Aggregation type.} Table~\ref{tab:agg} adds mean aggregation (3 seeds, medians; the sum and max rows use the same seeds 0--2 from the main group). Without steps, mean has the lowest S64 median ($0.180$) of the three, but with only three seeds, two of which fail for max, we do not claim a ranking. With steps, mean has a higher S64 median than sum ($0.728$ against $0.386$) and the highest S8 and S16 medians of all six rows, but it is the only step-supervised system with no failing seed at D64. One of the three mean-aggregation seeds without steps exceeds the failure threshold at D64. This ablation therefore shows that failures at 64 nodes are not confined to sum and max; it does not show that mean is better.

\begin{table}[t]\centering\caption{Aggregation ablation, median MAE over seeds 0--2; in parentheses the number of seeds with MAE $>1$ or non-finite.}\label{tab:agg}
\small\begin{tabular}{lcccc}\toprule
 Aggregation & S8 & S16 & S64 (fail) & D64 (fail) \\\midrule
sum & 0.016 & 0.026 & 0.417 (0/3) & 0.464 (0/3) \\
mean & 0.019 & 0.065 & 0.180 (0/3) & 0.363 (1/3) \\
max & 0.008 & 0.036 & 1.789 (2/3) & 2.450 (2/3) \\
sum + steps & 0.036 & 0.107 & 0.386 (0/3) & div. (3/3) \\
mean + steps & 0.078 & 0.203 & 0.728 (0/3) & 0.199 (0/3) \\
max + steps & 0.026 & 0.081 & 2.502 (2/3) & 4.709 (2/3) \\
\bottomrule\end{tabular}

\end{table}

\textbf{Test-time steps.} Table~\ref{tab:sweep} and Fig.~\ref{fig:steps} rerun the two step-supervised systems (seeds 0--2) with the number of message-passing steps at test time scaled to $0.5$, $1$, $2$ and $4$ times $n-1$ (training always used $n-1$; the $\times1$ column is the main-group result for the same seeds; $\times4$ was added after the registered sweep). The final-only systems were not swept, so nothing here applies to them. The median error increases with the multiplier in all four blocks: for max+steps at S64 it is $0.28$ at $\times0.5$ and $2.50$ at $\times1$, rising to $3\times10^{2}$ at $\times2$ and $10^{7}$ at $\times4$; for sum+steps it rises more slowly ($0.39$ at $\times1$, $2.73$ at $\times4$). Per seed the picture is less uniform. At S16 the error at $\times2$ and at $\times4$ is above that at $\times1$ for all six models. At S64 four of the six models get worse by a factor of more than eight at $\times4$, but max+steps seed 0 ($0.21$ to $0.23$) and sum+steps seed 1 ($0.46$ throughout) change little. Halving the steps lowers the error in every model and both cells. So for these six models extra rounds are not harmless: in most of them the learned update does not leave a finished computation unchanged, as the relaxation of Eq.~\eqref{eq:bf} would. We did not measure the number of rounds Bellman-Ford needs to converge on these graphs (usually fewer than $n-1$), so we cannot say whether $\times0.5$ is better because it avoids accumulated error or because fewer rounds suffice; both are possible.

\begin{table}[t]\centering\caption{Test-time step multiplier for the step-supervised systems: MAE per seed and median (seeds 0--2). Values of 100 or more to one significant figure. $\dagger$: not in the registered sweep (post hoc).}\label{tab:sweep}
\small\begin{tabular}{lllcccc}\toprule
 System & Cell & & $\times$0.5 & $\times$1 & $\times$2 & $\times$4$^\dagger$ \\\midrule
max+steps & S16 & median & 0.039 & 0.081 & 0.44 & 2.04 \\
 & & seed 0 & 0.039 & 0.057 & 0.097 & 0.12 \\
 & & seed 1 & 0.069 & 0.15 & 0.44 & 3.86 \\
 & & seed 2 & 0.031 & 0.081 & 0.59 & 2.04 \\
\midrule
max+steps & S64 & median & 0.28 & 2.50 & 3e2 & 1e7 \\
 & & seed 0 & 0.16 & 0.21 & 0.23 & 0.23 \\
 & & seed 1 & 0.28 & 5.15 & 7e4 & 8e12 \\
 & & seed 2 & 0.47 & 2.50 & 3e2 & 1e7 \\
\midrule
sum+steps & S16 & median & 0.067 & 0.11 & 0.19 & 0.35 \\
 & & seed 0 & 0.077 & 0.11 & 0.19 & 0.35 \\
 & & seed 1 & 0.067 & 0.12 & 0.61 & 0.80 \\
 & & seed 2 & 0.054 & 0.073 & 0.11 & 0.21 \\
\midrule
sum+steps & S64 & median & 0.31 & 0.39 & 0.46 & 2.73 \\
 & & seed 0 & 0.23 & 0.33 & 0.42 & 2.73 \\
 & & seed 1 & 0.45 & 0.46 & 0.46 & 0.46 \\
 & & seed 2 & 0.31 & 0.39 & 5.22 & 5e4 \\
\bottomrule\end{tabular}

\end{table}

\begin{figure}[t]\centering\includegraphics[width=\columnwidth]{steps_sweep.pdf}
\caption{Step-supervised systems: median MAE (dots: seeds 0--2) versus test-time steps divided by $n-1$, sparse family, 16 and 64 nodes (the two panels have different vertical ranges). All values are finite and none is clipped.}\label{fig:steps}\end{figure}
